另一条消息路径采用 Kafka。购买权威和消费者事务保持前述接口，完整 EV 由适配器写入 topic 分区，消费位移由相应 consumer group 保存。以下采用 Kafka 4.4.0、三 broker、三分区与每分区三副本，配置 \code{acks=all}、\code{min.insync.replicas=2}，并关闭消费者自动提交和自动保存待提交位移\footnote{自动 offset store 是相应客户端保存待提交位置的设置，与 broker 上的已提交位移分开。客户端选项名与默认值按所用实现核定，本章由业务事务成功后的路径显式推进。}。

\subsection{分区日志与副本保存}

topic 是日志的逻辑集合，每个分区拥有有序追加记录和独立 offset。offset 定位某分区内的一次日志位置，跨分区各有自己的序列。本章以 EV 作为 record key 参与分区，稳定选择同一事件的消息分区。topic 分区数量变化会影响通常的键映射，要求长期固定实体归属时，应用对分区变化另行设计路由与迁移。

日志采用有限保留时，消费组位移本身无法保存已经被删除的原消息。购买权威使用第~\ref{sec:ch16-handoff}~节的公共交接闭环：尚无可信消费回执的 EV 持续列入恢复扫描，必要时把同一原报文写到新的 offset。追加确认先持久写为 \code{published}，经核验的消费回执再共同推进 \code{handoff\_confirmed} 和履约已接受。Kafka 的 offset 记录本次传播位置，原 EV 和 F 连接同一事件与履约；消费组恢复、日志保留与业务交接因而可以各自推进。

\CSixteenFigure{ch16-10-kafka}{topic 分区保存日志位置，消费组各自保存读取进度，EV 继续标识业务事件}{fig:ch16-kafka}

每个分区有 leader 和 follower。ISR 即同步副本集合，维护当前达到相应同步条件的副本。生产者设置 \code{acks=all} 以后，leader 的追加确认依据 ISR 与高水位推进等条件完成；\code{min.insync.replicas=2} 约束本配置在同步副本不足时拒绝这类追加。ISR 可以随副本进度与故障变化，all 指向当前协议维护的同步集合~\cite{kafka44-replication}。

在锁定版本的实现中，leader 先检查当前 ISR 是否达到 minISR，追加以后，等待所需 offset 达到可确认水位，并再次考虑相应 ISR 条件。ISR 收缩前后可能产生不同错误，追加后错误保留结果未知的含义。重试由生产协议和业务原身份共同处理。为了保持已确认数据的故障保证，副本放置、同步条件、选主政策和磁盘保存都属于配置的一部分~\cite{kafka44-replication}。

Kafka 通常利用操作系统页缓存与复制获得吞吐，追加到副本日志和每副本逐条 fsync 是不同动作。\code{acks=all} 的确认规则依赖副本日志的保存，逐条强制刷盘则由另行配置的路径决定。进程退出、机器掉电以及多个相关副本同时丢失的条件分别检查，才能解释设计依赖哪份继续保存的日志~\cite{kafka44-design}。

KRaft 维护集群元数据的控制面共识，例如分区副本和 leader 等配置事实。topic 分区数据走上述 leader、ISR 和日志水位路径，购买业务权威又是另一组完整状态机。三者都涉及复制，所复制的对象、成员资格和提交规则各有位置~\cite{kafka44-design}。

\subsection{PID、epoch 与 sequence 怎样识别协议重试}

幂等生产者取得 producer ID，通常简称 PID。记录批次还携带生产者 epoch 和分区内 sequence。PID 定位生产协议实例，epoch 区分同一相关生产身份的代次，sequence 标记该生产者在特定分区的批次先后。broker 保存相应生产状态，用这些信息识别协议重试和不合顺序的批次~\cite{kafka44-producer}。

假设生产者把一批包含 \code{EV:b1:1} 的记录写入分区 0，broker 已经追加，\code{ProduceResponse} 却在网络中丢失。客户端重试该批次时沿用相应 PID、epoch 和 sequence。broker 从保存的批次元数据识别重复，返回原追加位置。锁定版本的日志路径先检查重复批次，命中以后回原 first/last offset，新的批次才进入追加与生产状态更新~\cite{kafka44-replication}。

epoch 较旧或 sequence 与已保存顺序不相容的来件会进入相应错误路径。客户端按该协议处理重试与恢复。生产状态有自己的保留和恢复机制，业务的长期原身份则保存在购买权威与消费数据库。

一个新的普通生产者进程可以取得新的 PID。它从 Outbox 再次发送同 EV 时，这是新生产协议身份下的一次应用发送，Kafka 能够把它追加到另一个 offset。\code{enable.idempotence=true} 保护相应生产会话内的协议重试，应用跨进程恢复继续依靠完整 EV 与消费者 I/B。图~\ref{fig:ch16-producer} 将这两种重试放到各自的身份范围。

\CSixteenFigure{ch16-20-producer}{原批次协议重试恢复原位置，新普通生产者再次发送同 EV 可以取得新位置}{fig:ch16-producer}

丢确认可以具体落在两个位置。原生 \code{ProduceResponse} 丢失以后，客户端在等待期限内报告超时，分区中已经存在一条原 EV。发布器重启后以新的普通生产者发送同 EV，再取得另一个 offset。日志中两条完整报文相同，消费者业务事务返回同一 F501。若在客户端已经收到成功回调以后进程退出，形成的则是确认到达与 Outbox 标记之间的窗口。

\code{transactional.id} 属于另一种稳定生产配置。事务应用用它标识跨重启恢复的事务生产者，事务协调者和 epoch 机制可以使旧实例失去继续提交的资格。选择 \code{transactional.id} 时，要对应一份明确的事务生产职责，避免两项独立职责误用同一名字互相围栏。它和 $S$、C、EV 分别定位业务意图、付款责任与事件~\cite{kafka44-producer}。

\subsection{consumer group 和活旧消费者}

consumer group 把多个进程组织成一份逻辑订阅。同一组当前把每个分区分配给一个成员，不同组可以独立读取同一日志并保存自己的位移。一个组里的消费者退出、新成员加入或分区集合变化，可以触发重平衡，协调者重新分配分区~\cite{kafka44-consumer}。

Kafka 区分本地 position 与 committed offset。\code{poll} 返回记录以后，本地下一读取位置已经推进；提交位移保存下一次恢复应从哪里开始。分区 0 的 offset 0 处理完成后，通常提交下一位置 1。实际日志还包含事务控制记录等对象，业务记录的 offset 可以有空隙，消费者按收到记录的真实位置推进~\cite{kafka44-consumer}。

消费组通过心跳和 poll 等活性条件判断成员。\code{session.timeout} 用于相关会话失联判断，\code{max.poll.interval} 限制应用长期不继续 poll。应用持续活着却卡在业务处理时，旧成员可以失去分区，后继成员从已提交位移取得相同记录。旧进程里的业务线程和已经发出的 HTTP 请求却可以继续运行~\cite{kafka44-consumer}。

\CSixteenFigure{ch16-11-rebalance}{组内重分配撤销位移提交资格，旧外部请求由消费域的原身份仲裁}{fig:ch16-rebalance}

图~\ref{fig:ch16-rebalance} 中 A 已经取到 EV 但长期不再 poll，B 随后取得分区 0 和同一 EV。A 与 B 都向消费者数据库发送原事件，两个 HTTP 调用的区间重叠。数据库的事务仲裁 I/B，两次返回同一 F501。A 随后提交旧组位移，协调者因成员资格失效拒绝该提交。位移拒交发生在 Kafka 域，数据库里的单一效果发生在另一个仲裁域。

这也给出消费者内部并发的要求。把处理交给线程池可以让 poll 继续保持活性，应用要暂停已交给处理线程但尚未完成的分区，并分别维护在途记录和连续完成水位。分区被撤销时，停止分派新工作，处理已有工作、等待或转交其持久责任，再按当前资格提交合法位置。外部效果按 I/B 和 dispatch 身份核验，后继消费者从原关联恢复未完成调用。

自动提交若先于业务事务完成，可能让后继成员从更高位置开始，漏掉尚未提交的效果。本章关闭自动提交与自动 offset store，由业务提交成功以后显式提交该记录的下一位置。若业务成功、位移提交响应丢失，恢复后再次读取原 EV，原 Inbox 返回原回执。

冲突与暂时失败也连接这两个位置。数据库返回 503 以后，本地 position 可能已经前移，应用回到原 offset 重试。确定异载荷则保存阻断记录并停住位移前缀。数据库成功、位移失败会带来重复读取，数据库失败、位移成功却会使责任越过未完成记录，应用的提交路径因此只接受已完成前缀。

\subsection{Kafka 事务的原子范围}

Kafka 事务把多个分区的输出追加和消费组位移更新放入同一提交决定。一个流处理程序读取输入 offset 0，计算输出，把输出记录和下一消费位移 1 送入同一事务。提交后，使用 \code{read\_committed} 的下游能读到已提交输出，消费组从 1 恢复。中止后，相应输出按事务状态过滤，组位移保持原值~\cite{kafka44-design}。

\CSixteenFigure{ch16-12-transactions}{Kafka 事务共同提交输出日志与消费位移，外部数据库和付款各有自己的提交}{fig:ch16-transactions}

\code{read\_committed} 读取到相应稳定位置，过滤中止事务的记录，尚未决定的事务也影响可见前缀。\code{read\_uncommitted} 可以看到后来被中止的事务数据。事务结果通过控制记录等元数据表达，所以 offset 0 与 2 之间缺少一条业务记录，可以来自这类协议记录~\cite{kafka44-consumer}。

事务中止以后，本地 consumer position 保留已读取进度。应用显式取得已提交位移并 \code{seek}，或重建消费者使它按已提交位置恢复，再重做原输入。只调用 \code{abortTransaction} 便继续 poll，可能跨过本次已经读过却没有成功提交的输入。内存计算状态和临时输出同样要恢复到对应输入边界~\cite{kafka44-consumer}。

事务生产配置要把事务身份、消费组元数据与当前分配共同带入流程。旧生产者被新实例围栏，保护的是相应事务生产职责；旧消费者在组内重平衡后提交位移被拒，保护的是消费资格。二者可以协作，分别有自己的参与者与错误路径。

通常所称 Kafka 的 EOS，即恰好一次处理语义，在这里对应从 Kafka 输入到 Kafka 输出及消费位移的完整事务链路。付款 HTTP 请求、消费者 SQL 事务和外部发货各自在自己的保存域提交。即使把 offset 写进 Kafka 事务，数据库已经提交以后应用退出，也会再次面对数据库效果。购买系统把这些外部域的原身份和持久责任保存下来，继续使用 I/B、原 C/T 和 dispatch 协议连接它们~\cite{kafka44-design}。

外部数据库也可以在自己的同一事务里保存输出与消费位置，这样数据库恢复时能够共同恢复效果与水位。采用这种设计时，数据库中的位置成为相应输出恢复的依据，分区重新分配与旧工作者写资格也在资源端处理。本章的参考路径使用完整 EV 的 Inbox 与业务守卫，Kafka 位移负责高效继续读取，业务原关系负责重复仲裁~\cite{kafka44-consumer}。

\subsection{消息子系统回到整体架构}

现在把消息的两个备选部署放回第一节的全景。购买分区拥有会话、$S$ 绑定、C/T、退款占额、原结果和 Outbox。RabbitMQ 路径把原 EV 放进队列并交付消费者，Kafka 路径把原 EV 追加到分区日志并维护各组读取位置。两条路径使用同一事件接口，消费事务核验相同的原事件与初次履约关联。一次部署选其中一条，消息域的配置和容量按选定路径落实。

图~\ref{fig:ch16-panorama} 按责任呈现完整运行链。同步受理路径只走入口、路由与所属购买组，返回固定 $A_J$。付款路径从 C/T 取得资格，沿原 scope/key 到平台，再由证据入口形成 $R_J$。传播路径从 Outbox 取得完整 EV，经过消息域到 I/B/F，持久回执返回购买权威后完成业务交接。发货路径从消费库意图到独立效果域。查询路径读购买权威和已应用的交接进度，组成原受理、原付款与当前进度。每条路径都能在自己的起点找到持续保存的工作。

\begingroup
\renewcommand{\CSixteenFigureWidth}{.90\linewidth}
\CSixteenFigure{ch16-21-panorama}{同步受理、付款、传播、发货与进度查询各有责任起点，恢复控制按这些起点枚举}{fig:ch16-panorama}
\endgroup

这种组织也决定故障后的恢复次序。购买组恢复多数并确认可读前缀以后，先恢复原绑定和资格，再扫描 C/T、TX、退款和 Outbox。平台已经完成而本地尚未形成 $R_J$ 的操作，通过原关联补齐结果，随后才出现可发布 EV。消息系统恢复后，发布器处理未获确认的事件，交接恢复者继续核查尚无回执的原 EV，消费者在原 I/B 恢复重复投递。履约数据库恢复后，dispatch 扫描器独立枚举未完成意图。各扫描器可以并行工作，但每项实际动作只消费其已经满足的持久前提。

控制面位于这些业务路径侧面。它改变归属和执行资格，业务数据组在本地检查激活或冻结状态。迁移器冻结完整可变图像、导入、交接并激活以后，原 C、EV 和 dispatch 关联保持原字面，后继工作者在新位置继续。路由缓存刷新解决寻找承载组的问题，原身份解决找到以后应恢复哪项业务的问题。分区快照还要包含会话关闭、退款占用和未完成责任，才能让新位置具有同一份继续依据。

运行观察同样沿责任分组。入口看新受理与恢复读取的速率，购买组看多数提交和未决 C 的年龄，平台适配器看查询窗口与未知期限，消息侧看未发布量、队列或日志积压，消费侧看阻断前缀、F 和未完成发货意图。某一指标超过预算时，控制器收紧对应的新责任准入，并给旧责任保留恢复能力。例如平台长时间不可用时限制新 C，消息恢复时留出高于持续到达率的消费余量，dispatch 变慢时限制并发而保留意图。

整条链随完整业务事实在新位置继续。副本切换恢复已确认的分区前缀，跨组迁移交接完整状态和唯一写资格，灾难回退先重建旧绑定及外部事实，再开放可能重用身份的新受理。第十七章沿稳定分区与当前承载组的关系，展开复制、资格与迁移的内部结构。

\par
\let\code\BeforeSixteenCode
\setlength{\emergencystretch}{\BeforeSixteenEmergency}
\tolerance=\BeforeSixteenTolerance\relax
